\documentclass[11pt]{article}

\usepackage{amsmath,amssymb}
\usepackage{physics}
\usepackage[utf8]{inputenc}
\usepackage[letterpaper,margin=1.2in]{geometry}
\usepackage[colorlinks=true,allcolors=blue]{hyperref}
\usepackage[makeroom]{cancel}
\usepackage{xcolor}
\usepackage{mdframed}
\usepackage{booktabs}

%% Colored boxes for key results and warnings
\newmdenv[backgroundcolor=blue!7,linecolor=blue!50,linewidth=1pt,
          innertopmargin=7pt,innerbottommargin=7pt,
          frametitle={\textbf{Key result}},
          frametitlebackgroundcolor=blue!15]{keyresult}

\newmdenv[backgroundcolor=orange!7,linecolor=orange!60,linewidth=1pt,
          innertopmargin=7pt,innerbottommargin=7pt,
          frametitle={\textbf{Common confusion -- read carefully}},
          frametitlebackgroundcolor=orange!15]{warning}

\newmdenv[backgroundcolor=green!7,linecolor=green!50,linewidth=1pt,
          innertopmargin=7pt,innerbottommargin=7pt,
          frametitle={\textbf{Physical interpretation}},
          frametitlebackgroundcolor=green!15]{physics}

%% Notation matching paper 2 notation.tex
\newcommand{\iunit}{\imath}
\newcommand{\Avec}{\mathbf{A}}
\newcommand{\Evec}{\mathbf{E}}
\newcommand{\Bvec}{\mathbf{B}}
\newcommand{\phat}{\hat{\mathbf{p}}}
\newcommand{\rhat}{\hat{\mathbf{r}}}
\newcommand{\Vnl}{V^{\mathrm{nl}}}
\newcommand{\HLG}{\hat{H}^{\mathrm{LG}}}
\newcommand{\HVG}{\hat{H}^{\mathrm{VG}}}
\newcommand{\VKS}{V_{\mathrm{KS}}}

\title{%
  Electromagnetic Perturbations in RT-TDDFT:\\
  From the Plane Wave to the Long-Wavelength Approximation\\[4pt]
  {\large Pedagogical Derivation --- Revised}%
}
\author{Working notes for RMG Paper 2 (Jacek Jakowski)}
\date{April 2026 --- Revision 2}

\begin{document}
\maketitle

\begin{abstract}
These notes derive the interaction Hamiltonian for RT-TDDFT starting
from a monochromatic electromagnetic plane wave.
The central organizing principle is a clean separation of two
physically distinct types of perturbation:
\textbf{Track~A} (electromagnetic field perturbations, enter through
the vector potential in the kinetic energy, require gauge treatment)
and \textbf{Track~B} (scalar potential perturbations, added directly
to the KS potential, no gauge involved).
The confusion between these two tracks --- which share the spatial
factor $e^{\iunit\mathbf{q}\cdot\mathbf{r}}$ but are otherwise
completely different --- is identified and resolved.
\end{abstract}

\tableofcontents
\newpage

%%=================================================================
\section{Conventions and Notation}
%%=================================================================

\noindent\textbf{Units.} Atomic units throughout unless stated:
$\hbar=m_e=e=4\pi\varepsilon_0=1$.
Speed of light $c=1/\alpha\approx137$ (kept explicit where it
clarifies sign conventions).
Electron charge $q_e=-1$ (i.e.\ charge magnitude $e=1$, electron
charge negative).

\noindent\textbf{Imaginary unit.} $\iunit$ (dotless $i$), consistent
with paper~1 and paper~2.

\noindent\textbf{Gauge.} We work in the Coulomb gauge
($\nabla\cdot\Avec=0$, scalar potential $\phi=0$ in vacuum) unless
a specific gauge is indicated.

%%=================================================================
\section{The Two Tracks: a Roadmap}
%%=================================================================

Before any derivation, we establish the central conceptual
distinction that organizes everything that follows.

\begin{warning}
The factor $e^{\iunit\mathbf{q}\cdot\mathbf{r}}$ appears in
\emph{two completely different contexts} in RT-TDDFT.
They share the same spatial structure but enter the Hamiltonian
in different places, probe different physics, and require different
theoretical treatment.
Conflating them is the source of significant confusion.
\end{warning}

\bigskip
\noindent\textbf{Track A --- Electromagnetic field perturbations.}
The plane wave $\Avec(\mathbf{r},t)
= A_0\hat{\varepsilon}\,e^{\iunit(\mathbf{q}\cdot\mathbf{r}-\omega t)}$
couples to electrons through \emph{minimal substitution}
$\phat\to\phat+\Avec$, entering the kinetic energy operator.
Because it appears in the kinetic energy, gauge freedom matters:
the length gauge (LG) and velocity gauge (VG) represent the same
physics differently, and the choice of gauge has practical
consequences for periodic systems.
The long-wavelength approximation ($\mathbf{q}\to0$) converts the
spatially varying $\Avec(\mathbf{r},t)$ to a spatially uniform
$\Avec(t)$ and gives the ``momentum kick'' perturbation used in
this paper.

\bigskip
\noindent\textbf{Track B --- Scalar potential perturbations.}
A time-dependent potential $\delta V(\mathbf{r},t)$ is added
directly to the Kohn-Sham effective potential $\VKS$,
with no vector potential and no modification of the kinetic energy.
No gauge transformation is needed or relevant; the Coulomb gauge is
maintained throughout.
The specific choices $\delta V = V_0 e^{\iunit\mathbf{q}\cdot\mathbf{r}}\delta(t)$
(density-wave kick) and
$\delta V = V_0/|\mathbf{R}-\mathbf{r}|\,\delta(t)$ (Coulomb-matrix
kick) both belong to this track.

\bigskip
\begin{center}
\begin{tabular}{p{2.4cm}p{3.6cm}p{3.6cm}p{3.0cm}}
\toprule
& \textbf{Track A (EM field)} & \textbf{Track B (scalar pot.)} & \\
\midrule
Perturbation
  & $\Avec(\mathbf{r},t)\cdot\phat$ in $\hat{T}$
  & $\delta V(\mathbf{r},t)$ added to $\VKS$
  & \\
Gauge relevant?
  & \textbf{Yes} --- LG vs VG
  & \textbf{No}
  & \\
$\hat{\mathbf{r}}$ problem?
  & Yes (LG for PBC)
  & No ($e^{\iunit\mathbf{q}\cdot\mathbf{r}}$ is periodic)
  & \\
Response probed
  & current--current: $\sigma(\omega)$
  & density--density: $\chi(\mathbf{q},\omega)$
  & \\
Observable
  & $\varepsilon(\omega)$, optical spectra
  & $S(\mathbf{q},\omega)$, EELS, IXS
  & \\
$\mathbf{q}$ constrained?
  & Yes: $|\mathbf{q}|=\omega/c$ (photon)
  & No: $\mathbf{q}$ free parameter
  & \\
This paper
  & \checkmark (long-wavelength, VG)
  & future work
  & \\
\bottomrule
\end{tabular}
\end{center}

%%=================================================================
\section{Track A: Electromagnetic Field Perturbations}
%%=================================================================

%%-------------------------------------------------------------------
\subsection{The plane wave vector potential}
%%-------------------------------------------------------------------

A monochromatic plane wave in the Coulomb gauge:
\begin{equation}
  \Avec(\mathbf{r},t)
  = A_0\,\hat{\varepsilon}\,
    e^{\iunit(\mathbf{q}\cdot\mathbf{r}-\omega t)}
    + \text{c.c.}
  \label{eq:A}
\end{equation}
with $|\mathbf{q}|=\omega/c$ and transversality
$\hat{\varepsilon}\cdot\mathbf{q}=0$.

The electric and magnetic fields:
\begin{align}
  \Evec(\mathbf{r},t)
  &= -\frac{\partial\Avec}{\partial t}
   = \iunit\omega A_0\,\hat{\varepsilon}\,
     e^{\iunit(\mathbf{q}\cdot\mathbf{r}-\omega t)}
     + \text{c.c.},
  \label{eq:E}
  \\
  \Bvec(\mathbf{r},t)
  &= \nabla\times\Avec
   = \iunit\mathbf{q}\times
     (A_0\hat{\varepsilon})\,
     e^{\iunit(\mathbf{q}\cdot\mathbf{r}-\omega t)}
     + \text{c.c.}
  \label{eq:B}
\end{align}

Derivation of eq.~(\ref{eq:B}): apply
$\nabla\times(f\mathbf{v})=(\nabla f)\times\mathbf{v}$ with
$f=e^{\iunit\mathbf{q}\cdot\mathbf{r}}$,
$\nabla f = \iunit\mathbf{q}\,e^{\iunit\mathbf{q}\cdot\mathbf{r}}$,
constant $\mathbf{v}=A_0\hat{\varepsilon}$.

\begin{physics}
All three fields $\Avec$, $\Evec$, $\Bvec$ carry the same spatial
phase $e^{\iunit\mathbf{q}\cdot\mathbf{r}}$.
Therefore when we make the long-wavelength approximation
$e^{\iunit\mathbf{q}\cdot\mathbf{r}}\approx1$, we simultaneously
make $\Avec$, $\Evec$, and $\Bvec$ spatially uniform.
The magnetic field does not disappear --- it becomes spatially
uniform.
But a spatially uniform $\Bvec$ corresponds to a spatially linear
$\Avec$ (symmetric gauge), not to zero field.
In the long-wavelength limit we are saying $|\mathbf{q}|\to0$,
which gives $\Bvec\to\mathbf{q}\times\hat{\varepsilon}\to0$.
So the magnetic field does vanish in the strict $\mathbf{q}=0$ limit.
\end{physics}

%%-------------------------------------------------------------------
\subsection{Minimal coupling and the interaction Hamiltonian}
%%-------------------------------------------------------------------

Minimal substitution $\phat\to\phat+\Avec(\mathbf{r},t)$ for an
electron (charge $-1$, giving $\phat\to\phat+\Avec$ in a.u.):
\begin{equation}
  \hat{H}(t)
  = \frac{1}{2}(\phat+\Avec)^2 + \VKS
  = \underbrace{-\frac{\nabla^2}{2}+\VKS}_{\hat{H}_0}
  + \underbrace{\Avec\cdot\phat}_{\hat{H}^{(1)}}
  + \underbrace{\frac{1}{2}\Avec^2}_{\hat{H}^{(2)}}.
  \label{eq:H-expand}
\end{equation}
In linear response, drop $\hat{H}^{(2)}\sim A_0^2$.
The electromagnetic interaction Hamiltonian is:
\begin{equation}
  \hat{H}_\mathrm{int}(\mathbf{r},t)
  = \Avec(\mathbf{r},t)\cdot\phat
  = A_0\,\hat{\varepsilon}\cdot\phat\;
    e^{\iunit(\mathbf{q}\cdot\mathbf{r}-\omega t)}
    + \text{c.c.}
  \label{eq:Hint}
\end{equation}

\begin{warning}
The factor $e^{\iunit\mathbf{q}\cdot\mathbf{r}}$ in
eq.~(\ref{eq:Hint}) multiplies the \emph{momentum operator}
$\hat{\varepsilon}\cdot\phat$ --- it enters through the kinetic
energy, not through a scalar potential.
This is Track~A.
Track~B perturbations of the form $V_0 e^{\iunit\mathbf{q}\cdot\mathbf{r}}$
are added to $\VKS$, not to the kinetic energy.
They look the same spatially but are physically different.
\end{warning}

%%-------------------------------------------------------------------
\subsection{The multipole expansion}
%%-------------------------------------------------------------------

Taylor expand $e^{\iunit\mathbf{q}\cdot\mathbf{r}}$ in the matrix
element of $\hat{H}_\mathrm{int}$:
\begin{equation}
  e^{\iunit\mathbf{q}\cdot\mathbf{r}}
  = 1
  + \iunit(\mathbf{q}\cdot\rhat)
  - \tfrac{1}{2}(\mathbf{q}\cdot\rhat)^2
  + \cdots
  \label{eq:Taylor}
\end{equation}
\begin{equation}
  \langle\psi_b|\hat{H}_\mathrm{int}|\psi_a\rangle
  \propto
  \underbrace{
    \langle\psi_b|\hat{\varepsilon}\cdot\phat|\psi_a\rangle
  }_{\substack{\text{order 0}\\[2pt]\text{E1 velocity form}}}
  +\;\iunit
  \underbrace{
    \langle\psi_b|(\mathbf{q}\cdot\rhat)(\hat{\varepsilon}\cdot\phat)
    |\psi_a\rangle
  }_{\substack{\text{order 1}\\[2pt]\text{M1 + E2}}}
  + \cdots
  \label{eq:ME-expand}
\end{equation}

\noindent Relative size of corrections for optical light
($\lambda=400$~nm, $a=10$~\AA):
\begin{equation}
  |\mathbf{q}|\,a
  = \frac{2\pi}{\lambda}\cdot a
  = \frac{2\pi}{7559\,a_0}\cdot18.9\,a_0
  \approx 0.016.
\end{equation}
The M1+E2 term is $\sim$1.6\% of E1.
The long-wavelength approximation is excellent.

%%-------------------------------------------------------------------
\subsection{Long-wavelength approximation: A(r,t) becomes A(t)}
%%-------------------------------------------------------------------

Setting $e^{\iunit\mathbf{q}\cdot\mathbf{r}}\approx1$:
\begin{align}
  \Avec(\mathbf{r},t)&\;\longrightarrow\;\Avec(t)
  = A_0\hat{\varepsilon}\,e^{-\iunit\omega t}+\text{c.c.}
  \label{eq:A-uniform}
  \\
  \hat{H}_\mathrm{int}(t) &= \Avec(t)\cdot\phat.
  \label{eq:Hint-LWA}
\end{align}
The vector potential is now spatially uniform (depends only on
time).
Crucially, a spatially uniform $\Avec(t)$ has:
\begin{equation}
  \Bvec = \nabla\times\Avec(t) = \mathbf{0}.
\end{equation}
The magnetic field is exactly zero.
The perturbation is purely electric, and no position operator
$\rhat$ appears in $\hat{H}_\mathrm{int}$ --- which is the key
advantage of the velocity gauge (Section~\ref{sec:VG}).

%%-------------------------------------------------------------------
\subsection{Crystal momentum conservation}
%%-------------------------------------------------------------------

For Bloch states $\psi_{n\mathbf{k}}
= e^{\iunit\mathbf{k}\cdot\mathbf{r}}u_{n\mathbf{k}}$,
the matrix element of the full EM interaction~(\ref{eq:Hint})
gives:
\begin{equation}
  \langle\psi_{m\mathbf{k}'}|
  e^{\iunit\mathbf{q}\cdot\mathbf{r}}\hat{\varepsilon}\cdot\phat
  |\psi_{n\mathbf{k}}\rangle
  \propto \delta_{\mathbf{k}',\mathbf{k}+\mathbf{q}}.
  \label{eq:k-cons}
\end{equation}
The photon delivers momentum $\mathbf{q}$ to the electron.
For optical light $|\mathbf{q}|\ll\pi/a$, so
$\Delta\mathbf{k}=\mathbf{q}\approx0$: all transitions are
\textbf{vertical} in the Brillouin zone.

%%-------------------------------------------------------------------
\subsection{Gauge transformation: length to velocity gauge}
\label{sec:VG}
%%-------------------------------------------------------------------

\subsubsection{Why the gauge transformation is needed --- and what
  it does}

The motivation for the LG$\to$VG gauge transformation is specific:
\textbf{to eliminate the position operator $\rhat$ from the
interaction Hamiltonian}.

In the length gauge, the long-wavelength interaction is
$\hat{H}_\mathrm{int}^\mathrm{LG}=\rhat\cdot\Evec(t)$, which
requires matrix elements
$\langle\psi_a|\hat{r}_\alpha|\psi_b\rangle$.
For periodic systems, $\hat{r}_\alpha$ is ill-defined because Bloch
states are extended over all space and $\hat{r}_\alpha$ is not
compatible with periodic boundary conditions.

The gauge transformation chosen --- the G\"{o}ppert-Mayer
transformation with gauge function $\chi=-\rhat\cdot\Avec^\mathrm{VG}$
--- is designed so that its time-derivative contribution
$-\iunit\hat{U}\partial_t\hat{U}^\dagger = +\rhat\cdot\Evec$
\emph{exactly cancels} the LG coupling.
The vector potential $\Avec$ that appears in the VG kinetic energy
$\frac{1}{2}(\phat+\Avec)^2$ is the replacement.

\begin{keyresult}
The gauge transformation LG$\to$VG is motivated by, and only by,
the need to eliminate $\hat{\mathbf{r}}\cdot\mathbf{E}$ from the
interaction.
If the perturbation does not contain $\hat{\mathbf{r}}$, there
is nothing to eliminate and the motivation for the gauge
transformation does not apply.
This is why Track~B (scalar potential) perturbations of the form
$V_0 e^{\iunit\mathbf{q}\cdot\mathbf{r}}\delta(t)$ require no
gauge transformation: $e^{\iunit\mathbf{q}\cdot\mathbf{r}}$ is
well-defined for periodic systems for any $\mathbf{q}\neq0$,
and it never appears in the kinetic energy.
\end{keyresult}

\subsubsection{The G\"{o}ppert-Mayer transformation: step by step}

Gauge potentials (LG and VG):
\begin{alignat}{2}
  &\text{LG:}\quad
  &&\Avec^\mathrm{LG}=\mathbf{0},\quad
  \phi^\mathrm{LG}=-\rhat\cdot\Evec(t);
  \label{eq:LG}
  \\
  &\text{VG:}\quad
  &&\Avec^\mathrm{VG}(t)=-\!\int_{-\infty}^t\Evec(t')dt',\quad
  \phi^\mathrm{VG}=0.
  \label{eq:VG}
\end{alignat}
Gauge function connecting them:
\begin{equation}
  \chi(\mathbf{r},t)
  = -\rhat\cdot\Avec^\mathrm{VG}(t).
  \label{eq:chi}
\end{equation}
Transformation rules:
\begin{align}
  \Avec &\to \Avec+\nabla\chi,
  &
  \phi &\to \phi-\frac{\partial\chi}{\partial t},
  \label{eq:gauge-rules}
  \\
  \Psi &\to \Psi' = \hat{U}\Psi,
  &
  \hat{U} &= e^{-\iunit\chi}
           = e^{\iunit\rhat\cdot\Avec(t)}.
  \label{eq:U-def}
\end{align}
The Hamiltonian transforms as:
\begin{equation}
  \hat{H}
  \to\hat{H}'
  = \hat{U}\hat{H}\hat{U}^\dagger
    - \iunit\hat{U}\frac{\partial\hat{U}^\dagger}{\partial t}.
  \label{eq:H-transform}
\end{equation}

Starting from $\HLG = -\nabla^2/2 + \VKS + \Vnl + \rhat\cdot\Evec$,
we evaluate four contributions.

\paragraph{Step 1 --- Transform the kinetic energy.}
\begin{align}
  \hat{U}\hat{p}_\alpha\hat{U}^\dagger\,g
  &= e^{\iunit\mathbf{r}\cdot\Avec}
     (-\iunit\nabla_\alpha)
     (e^{-\iunit\mathbf{r}\cdot\Avec}g)
  \nonumber\\
  &= e^{\iunit\mathbf{r}\cdot\Avec}(-\iunit)
     \bigl[(-\iunit A_\alpha)e^{-\iunit\mathbf{r}\cdot\Avec}g
     + e^{-\iunit\mathbf{r}\cdot\Avec}\nabla_\alpha g\bigr]
  \nonumber\\
  &= (\hat{p}_\alpha - A_\alpha)\,g.
  \label{eq:U-p-step}
\end{align}
So $\hat{U}\phat\hat{U}^\dagger = \phat-\Avec$, giving:
\begin{equation}
  \hat{U}\!\left(-\frac{\nabla^2}{2}\right)\hat{U}^\dagger
  = \frac{1}{2}(\phat-\Avec)^2
  \;\xrightarrow{\text{electron charge }=-1}\;
  \frac{1}{2}(\phat+\Avec)^2.
  \label{eq:T-VG}
\end{equation}
The sign reversal $-\Avec\to+\Avec$ comes from the electron charge
$-1$ (in atomic units, the coupling is $q_e\Avec=-\Avec$, so the
sign cancels).

\paragraph{Step 2 --- Time-derivative term.}
\begin{align}
  -\iunit\hat{U}\frac{\partial\hat{U}^\dagger}{\partial t}
  &= -\iunit\,e^{\iunit\mathbf{r}\cdot\Avec}
     \frac{\partial}{\partial t}
     e^{-\iunit\mathbf{r}\cdot\Avec}
   = -\iunit(-\iunit\mathbf{r}\cdot\dot\Avec)
   = -\mathbf{r}\cdot\dot\Avec(t)
   = +\mathbf{r}\cdot\Evec(t).
  \label{eq:dt-step}
\end{align}
This is $+\rhat\cdot\Evec$, which \emph{exactly cancels} the LG
coupling $+\rhat\cdot\Evec$ in $\HLG$:
\begin{equation}
  \underbrace{+\rhat\cdot\Evec}_{\text{from }\hat{U}\HLG\hat{U}^\dagger}
  +\underbrace{(-\rhat\cdot\Evec)}_{\text{Step 2}}
  \xrightarrow{\text{cancel}}
  0.
  \label{eq:cancel}
\end{equation}

\paragraph{Step 3 --- Local potential.}
$\VKS(\mathbf{r})$ commutes with $\hat{U}(\mathbf{r},t)$:
$\hat{U}\VKS\hat{U}^\dagger = \VKS$.

\paragraph{Step 4 --- Nonlocal pseudopotential.}
$\Vnl$ is not diagonal in position space, so:
\begin{equation}
  \hat{U}\Vnl\hat{U}^\dagger
  = e^{\iunit\rhat\cdot\Avec}\Vnl e^{-\iunit\rhat\cdot\Avec}
  \approx \Vnl + \iunit\Avec(t)\cdot[\rhat,\Vnl]
  + \mathcal{O}(A_0^2).
  \label{eq:Vnl-step}
\end{equation}
The first-order correction $\iunit\Avec\cdot[\rhat,\Vnl]$ is the
Starace commutator, identical to the nonlocal correction in the
velocity operator (Section 2.5 of manuscript).

\paragraph{Collecting all steps:}
\begin{equation}
  \boxed{
  \HVG(t)
  = \frac{1}{2}(\phat+\Avec)^2
  + \VKS(\mathbf{r},t)
  + \Vnl
  + \iunit\Avec(t)\cdot[\rhat,\Vnl],
  }
  \label{eq:HVG}
\end{equation}
which contains no $\rhat$ in the interaction.
For Bloch states, $-\iunit\nabla\to-\iunit\nabla+\mathbf{k}$
from the Bloch factor, giving the $\mathbf{k}$-dependent form
used in the code.

\subsubsection{The delta-kick limit}

With $\Avec(t)=A_0\hat{\varepsilon}\,\delta(t)$:
\begin{itemize}
  \item At $t=0^+$: each orbital acquires phase
    $(1+\iunit A_0\hat{\varepsilon}\cdot\phat)\psi_{n\mathbf{k}}^{(0)}$
    (for small $|A_0|$).
  \item For $t>0$: $\Avec=0$, system evolves under
    $\hat{H}_0+\VKS[\rho(t)]$ alone.
  \item The current $J_\alpha(t>0)$ encodes the full optical
    response via FFT.
\end{itemize}

%%=================================================================
\section{Track B: Scalar Potential Perturbations}
%%=================================================================

%%-------------------------------------------------------------------
\subsection{Definition and gauge status}
%%-------------------------------------------------------------------

A scalar potential perturbation is an addition to $\VKS$:
\begin{equation}
  \hat{H}(t)
  = \hat{H}_0 + \delta V(\mathbf{r},t),
  \qquad
  \delta V \text{ added to } \VKS \text{ directly.}
  \label{eq:Track-B}
\end{equation}
There is no vector potential.
No minimal substitution is involved.
No gauge transformation is needed or motivated.
The Coulomb gauge ($\nabla\cdot\Avec=0$, $\phi=0$) is maintained
throughout by construction.

The electromagnetic field does not appear.
Track~B perturbations do not represent photons at all --- they
represent external probe potentials (from electron beams, neutrons,
or simply as mathematical tools to compute response functions).

%%-------------------------------------------------------------------
\subsection{The density-wave kick}
%%-------------------------------------------------------------------

\begin{equation}
  \delta V(\mathbf{r},t)
  = V_0\,e^{\iunit\mathbf{q}\cdot\mathbf{r}}\,\delta(t),
  \label{eq:density-wave}
\end{equation}
where $\mathbf{q}$ is a \emph{free parameter}, not constrained by
any dispersion relation.
This perturbation excites a density wave at wavevector $\mathbf{q}$
and all frequencies simultaneously (the delta function in time
is broadband in frequency).

The density response:
\begin{equation}
  \delta\rho(\mathbf{q},t)
  = \int_\Omega\delta\rho(\mathbf{r},t)\,
    e^{-\iunit\mathbf{q}\cdot\mathbf{r}}\,d^3r.
  \label{eq:delta-rho}
\end{equation}
After Fourier transform in time:
\begin{equation}
  \chi(\mathbf{q},\omega)
  = \frac{\widetilde{\delta\rho}(\mathbf{q},\omega)}{V_0},
  \qquad
  S(\mathbf{q},\omega) \propto \mathrm{Im}\,\chi(\mathbf{q},\omega).
  \label{eq:chi-S}
\end{equation}

\begin{physics}
The $e^{\iunit\mathbf{q}\cdot\mathbf{r}}$ factor in the density-wave
kick superficially resembles the EM plane wave factor in Track~A.
The crucial difference:
\begin{itemize}
  \item In Track~A: $e^{\iunit\mathbf{q}\cdot\mathbf{r}}$ multiplies
    $\phat$ in the kinetic energy; $\mathbf{q}$ is the photon
    wavevector, constrained to $|\mathbf{q}|=\omega/c$.
  \item In Track~B: $e^{\iunit\mathbf{q}\cdot\mathbf{r}}$ is the
    spatial profile of a scalar potential in $\VKS$; $\mathbf{q}$
    is a freely chosen parameter with no relation to $\omega$.
\end{itemize}
The Track~B perturbation with $\mathbf{q}=0$ gives a spatially
uniform scalar kick $V_0\delta(t)$, which shifts all orbital energies
equally and produces no response.
The Track~A perturbation with $\mathbf{q}=0$ gives the optical
momentum kick that we actually use.
\end{physics}

%%-------------------------------------------------------------------
\subsection{The Coulomb-matrix kick}
%%-------------------------------------------------------------------

\begin{equation}
  \delta V_\mathrm{Coul}(\mathbf{r},t)
  = \frac{V_0}{|\mathbf{R}-\mathbf{r}|}\,\delta(t),
  \label{eq:Coulomb-kick}
\end{equation}
with matrix elements
$\langle\phi_i|1/|\mathbf{R}-\rhat||\phi_j\rangle$.
This mimics a focused electron beam at position $\mathbf{R}$
[Jakowski et al.\ JCTC 2019].

Connection to the density-wave kick via Fourier expansion of the
Coulomb kernel:
\begin{equation}
  \frac{1}{|\mathbf{R}-\mathbf{r}|}
  = \frac{4\pi}{\Omega}
    \sum_{\mathbf{q}\neq0}
    \frac{e^{\iunit\mathbf{q}\cdot(\mathbf{r}-\mathbf{R})}}{|\mathbf{q}|^2}.
  \label{eq:Coulomb-Fourier}
\end{equation}
Comparing eqs.~(\ref{eq:density-wave}) and (\ref{eq:Coulomb-Fourier}):

\begin{keyresult}
The Coulomb kick at position $\mathbf{R}$ is a superposition of
all density-wave kicks $e^{\iunit\mathbf{q}\cdot\mathbf{r}}$
weighted by $e^{-\iunit\mathbf{q}\cdot\mathbf{R}}/|\mathbf{q}|^2$.
\begin{itemize}
  \item The $1/|\mathbf{q}|^2$ weight emphasizes small-$\mathbf{q}$
    (long-wavelength) modes, reflecting the long-range Coulomb
    interaction.
  \item Density-wave kicks: momentum-resolved response
    $\chi(\mathbf{q},\omega)$ --- ideal for plasmon dispersion,
    comparison with EELS/IXS.
  \item Coulomb kicks: real-space resolved response
    $\chi(\mathbf{R},\omega)$ --- ideal for STEM-EELS, spatially
    resolved spectroscopy.
  \item Both are Track~B: no gauge involved, $\delta V$ added
    directly to $\VKS$.
\end{itemize}
\end{keyresult}

%%-------------------------------------------------------------------
\subsection{Why $e^{\iunit\mathbf{q}\cdot\mathbf{r}}$ does not
  require a gauge transformation for $\mathbf{q}\neq0$}
%%-------------------------------------------------------------------

The position operator $\rhat$ is ill-defined for periodic systems
because $\langle\psi_{n\mathbf{k}}|\hat{r}_\alpha|\psi_{m\mathbf{k}}\rangle$
is not translationally invariant.
But $e^{\iunit\mathbf{q}\cdot\mathbf{r}}$ is a \emph{periodic
operator} for any $\mathbf{q}$ commensurate with the reciprocal
lattice.
Its matrix elements between Bloch states,
\begin{equation}
  \langle\psi_{m\mathbf{k}}|e^{\iunit\mathbf{q}\cdot\rhat}|\psi_{n\mathbf{k}}\rangle
  = \int_\Omega u_{m\mathbf{k}}^*\,e^{\iunit\mathbf{q}\cdot\mathbf{r}}\,
    u_{n\mathbf{k}}\,d^3r,
  \label{eq:periodic-me}
\end{equation}
are perfectly well-defined: $e^{\iunit\mathbf{q}\cdot\mathbf{r}}$
maps a Bloch state at $\mathbf{k}$ to one at $\mathbf{k}+\mathbf{q}$,
which is also a valid Bloch state.
There is no periodicity problem and no gauge transformation needed.

The LG difficulty ($\hat{r}_\alpha$ ill-defined for PBC) is
specifically a $\mathbf{q}=0$ problem.
Taking the limit $\mathbf{q}\to0$ in
$e^{\iunit\mathbf{q}\cdot\mathbf{r}}\approx1+\iunit\mathbf{q}\cdot\mathbf{r}$
recovers the ill-defined $\hat{r}_\alpha$ --- which is why we need
the gauge transformation in Track~A at $\mathbf{q}=0$.
At finite $\mathbf{q}$, no such problem arises.

%%=================================================================
\section{The Spin-Polarized Extension of Track B}
%%=================================================================

With collinear spin-polarized RT-TDDFT, Track~B perturbations can
be applied to the two spin channels independently to probe
different types of response:

\begin{center}
\begin{tabular}{p{3.5cm}p{3.5cm}p{4cm}p{2.0cm}}
\toprule
\textbf{Perturbation} & \textbf{Applied to} & \textbf{Observable}
  & \textbf{Experiment} \\
\midrule
$+V_0\,e^{\iunit\mathbf{q}\cdot\mathbf{r}}\delta(t)$ (both spins)
  & $\VKS^\uparrow$ and $\VKS^\downarrow$
  & Charge response $\chi_{\rho\rho}(\mathbf{q},\omega)$, $S(\mathbf{q},\omega)$
  & EELS, IXS \\
\midrule
$+V_0\,e^{\iunit\mathbf{q}\cdot\mathbf{r}}\delta(t)$ (spin $\uparrow$) /
$-V_0\,e^{\iunit\mathbf{q}\cdot\mathbf{r}}\delta(t)$ (spin $\downarrow$)
  & $\VKS^\uparrow$ and $\VKS^\downarrow$ oppositely
  & Longitudinal spin response $\chi_{m_z m_z}(\mathbf{q},\omega)$
  & Neutron (longitudinal channel) \\
\bottomrule
\end{tabular}
\end{center}

Note: transverse spin fluctuations (magnon dispersion) require
non-collinear spin dynamics (two-component spinors) and are beyond
the scope of the collinear implementation.

%%=================================================================
\section{Summary: Hierarchy of Approximations and Perturbations}
%%=================================================================

\begin{center}
\begin{tabular}{p{2.0cm}p{4.2cm}p{3.0cm}p{3.2cm}}
\toprule
\textbf{Track} & \textbf{Perturbation} & \textbf{Response}
  & \textbf{Experiment} \\
\midrule
A (EM, VG, this paper)
  & $\Avec(t)\cdot\phat\,\delta(t)$, uniform
  & $\sigma(\omega)$, $\varepsilon(\omega)$
  & Optical/UV spectroscopy \\
\midrule
A (EM, future)
  & $\Avec(\mathbf{r})\cdot\phat$, sym.\ gauge
  & ECD, MCD
  & Chiroptical spectroscopy \\
\midrule
B (scalar, future)
  & $V_0\,e^{\iunit\mathbf{q}\cdot\mathbf{r}}\delta(t)$
    in $\VKS$ (both spins)
  & $\chi_{\rho\rho}(\mathbf{q},\omega)$, $S(\mathbf{q},\omega)$
  & EELS, IXS, neutron (charge) \\
\midrule
B (scalar, future)
  & $\pm V_0\,e^{\iunit\mathbf{q}\cdot\mathbf{r}}\delta(t)$
    in $\VKS^{\uparrow/\downarrow}$
  & $\chi_{m_z m_z}(\mathbf{q},\omega)$
  & Neutron (longitudinal spin) \\
\midrule
B (scalar, Ref.\ [2019])
  & $V_0/|\mathbf{R}-\mathbf{r}|\,\delta(t)$
    in $\VKS$
  & $\chi(\mathbf{R},\omega)$
  & STEM-EELS \\
\bottomrule
\end{tabular}
\end{center}

\bigskip
\noindent\textbf{The single most important takeaway:}

\begin{keyresult}
The gauge transformation LG$\to$VG is needed in Track~A because
the long-wavelength ($\mathbf{q}\to0$) EM interaction introduces
$\rhat$, which is ill-defined for periodic systems.
Track~B perturbations never introduce $\rhat$ and require no gauge
transformation.
The shared factor $e^{\iunit\mathbf{q}\cdot\mathbf{r}}$ in Track~A
(EM wave) and Track~B (density-wave kick) is a superficial
resemblance; the two perturbations enter the Hamiltonian in
different places, probe different response functions, and connect
to different experimental observables.
\end{keyresult}

\end{document}
